% META: {"id": "1SPE-SUITES-CO-022", "chapitre": "1SPE-SUITES", "type_objet": "corrige", "exercice_ref": "1SPE-SUITES-EX-022", "capacites_codes": ["C1", "C2"], "sources_inspiration": [], "status": "generated"}
% BEGIN-VERIFY
% from sympy import symbols, simplify
% n = symbols('n', integer=True, nonnegative=True)
% u0 = 7
% r = 5
% u = u0 + r * n
% assert u.subs(n, 0) == 7
% assert u.subs(n, 1) == 12
% assert u.subs(n, 2) == 17
% assert u.subs(n, 3) == 22
% diff = u.subs(n, n+1) - u
% assert simplify(diff) == 5
% assert u.subs(n, 10) == 57
% # Plus petit n tel que u_n > 100 : 7 + 5n > 100 => n > 93/5 = 18.6 => n = 19
% assert u.subs(n, 18) == 97
% assert u.subs(n, 19) == 102
% END-VERIFY
\begin{corrige}{1SPE-SUITES-EX-022}

\textbf{Question 1 — Calcul des premiers termes et conjecture.}

On utilise la relation de récurrence $u_{n+1} = u_n + 5$.

\[
  u_1 = u_0 + 5 = 7 + 5 = 12.
\]
\[
  u_2 = u_1 + 5 = 12 + 5 = 17.
\]
\[
  u_3 = u_2 + 5 = 17 + 5 = 22.
\]

\commentaireMarge{Différence constante entre termes consécutifs : $+5$ à chaque étape.}

La différence entre deux termes consécutifs est constante et égale à $5$. On conjecture que $(u_n)$ est une suite arithmétique de raison $5$.

\medskip
\textbf{Question 2 — Démonstration de la nature arithmétique.}

Pour tout entier naturel $n$, on calcule :
\[
  u_{n+1} - u_n = (u_n + 5) - u_n = 5.
\]

\commentaireMarge{Définition d'une suite arithmétique : différence $u_{n+1} - u_n$ constante.}

La différence $u_{n+1} - u_n$ est constante et égale à $5$ pour tout $n \in \mathbb{N}$.

\textbf{Conclusion :} $(u_n)$ est une suite arithmétique de raison $r = 5$ et de premier terme $u_0 = 7$.

\medskip
\textbf{Question 3 — Expression de $u_n$ en fonction de $n$.}

Par la formule du terme général d'une suite arithmétique :
\[
  u_n = u_0 + n \times r = 7 + 5n.
\]

\commentaireMarge{Formule : $u_n = u_0 + nr$ pour une suite arithmétique.}

Ainsi, pour tout $n \in \mathbb{N}$ : $\boxed{u_n = 5n + 7}$.

\medskip
\textbf{Question 4 — Calcul de $u_{10}$ et détermination du premier rang tel que $u_n > 100$.}

\textit{Calcul de $u_{10}$ :}
\[
  u_{10} = 5 \times 10 + 7 = 50 + 7 = 57.
\]

\textit{Plus petit entier $n$ tel que $u_n > 100$ :}

On résout l'inéquation $u_n > 100$ :
\[
  5n + 7 > 100 \iff 5n > 93 \iff n > \frac{93}{5} = 18{,}6.
\]

\commentaireMarge{$n$ est un entier, donc on prend le plus petit entier strictement supérieur à $18{,}6$.}

Le plus petit entier vérifiant $n > 18{,}6$ est $n = 19$.

\textit{Vérification :}
\[
  u_{18} = 5 \times 18 + 7 = 97 \leq 100, \quad u_{19} = 5 \times 19 + 7 = 102 > 100. \checkmark
\]

Le plus petit entier $n$ tel que $u_n > 100$ est $\boxed{n = 19}$.

\baremeIndicatif{Q1 : 1 pt (calculer $u_1, u_2, u_3$ et conjecturer) ; Q2 : 3 pts (raisonner, communiquer : démonstration complète avec la définition) ; Q3 : 2 pts (calculer : formule du terme général) ; Q4 : 2 pts (calculer $u_{10}$ : 1 pt ; résoudre l'inéquation et conclure : 1 pt)}
\end{corrige}
